\section{引言}
\label{sec:introduction}

量子计算利用量子态的叠加、干涉与纠缠处理信息，为量子系统模拟及若干计算密集型问题提供了不同于经典计算的求解途径\cite{preskill2018quantum}。扩大可执行电路规模，需要增加量子比特数量、提高操作精度，并通过硬件控制与编译优化减少执行中的误差累积。中性原子处理器利用光镊阵列约束原子，以Rydberg相互作用实现纠缠门，并通过相干输运重构原子间的连接。数百量子比特阵列、高保真纠缠门和相干原子输运的实验进展\cite{bluvstein2024logicalprocessor,evered2023highfidelity,bluvstein2022coherenttransport}，使中性原子成为研究大规模量子计算的重要平台。

中性原子编译器通过安排原子位置、移动轨迹和激光脉冲，将逻辑电路转化为物理操作序列。分区中性原子架构将原子存储与纠缠门执行分配到不同区域，以存储隔离抑制空闲原子的Rydberg受激。跨区输运引入阱转移损耗并消耗相干时间，层间放置与输运安排直接影响电路保真度和执行时延\cite{stade2024abstractmodel,lin2025zac}。

层边界处的原子位置会改变后续门层的执行条件。暂时空闲的原子若保留在纠缠区，可省去往返存储区的阱转移，但也会经历中间层的Rydberg脉冲并产生额外受激。若将原子送回存储区，所选阱位成为后续重入的起点，并可能占用其他原子的输运通道。移动光阱的结构约束决定哪些轨迹可以并行，一个重排批次的时延由其中最长轨迹决定\cite{lin2025zac,stade2025routingaware}。门位、驻留和存储位共同影响当前层及后续层的物理代价。

Lin等人提出的ZAC在初始布局中利用全电路交互，在动态放置中通过相邻门层匹配复用纠缠区原子，并将后续双比特门的配对关系及其距离代价用于门位和存储位选择\cite{lin2025zac}。Stade等人的路由感知放置把轨迹兼容性与批次时延纳入A*代价\cite{stade2025routingaware}。评价非相邻门层之间的驻留选择，需要跟踪原子在中间层的位置及其受到的脉冲，确定当前驻留与回迁对后续受激损耗、阱转移和可行轨迹的影响。

本文提出层边界联合优化方法GA-LK\footnote{工程代码及实验配置见\url{https://github.com/zhouzixiang1/zac-ga}。}。该方法采用遗传搜索联合选择门位与驻留状态，并通过单射匹配为回迁原子分配存储位。各候选以可执行输运计算当前损失，再从执行后的位置、占位和时钟出发，预测后续多个门层的输运、空闲受激与相干损失。联合搜索处理当前选择之间的约束，多层前瞻则使这些选择兼顾非相邻门层的交互需求。

本文在ZAC论文采用的18个电路以及MQT QMAP示例集的154个输入电路上进行评估。相较于Lin等人的ZAC编译器\cite{lin2025zac}和Stade等人的路由感知放置方法\cite{stade2025routingaware}，GA-LK在两个基准集上均获得更高的模型保真度几何均值。
